%% kapa: exam=1 round=19 year=2008 period=2 form=A subject=economics subject_ko=경제원론 seq=6 number=6
%% source: Q-Net 감정평가사 기출문제 (cst00309) artlSeq=1001962 file=extracted/1008/1001962/2008년 제19회 감정평가사 1차시험 기출문제/2교시.hwp
%% conversion: kordoc 4.24.51 -> tools/gen1.py (auto); review=proofread-sonnet; flags=none
\begin{question}{6}
두 나라 사이에 교역이 이루어지는 기본 원리와 관련하여 옳은 설명을 모두 고르면?
\begin{framedtext}
ㄱ. 각국은 기회비용이 작은 재화를 생산한다. \par
ㄴ. 한 나라가 모든 재화에 절대적인 우위가 있는 경우 교역은 이루어지지 않는다. \par
ㄷ. 교역이 이루어지는 경우 한 나라가 이득을 보면 다른 나라는 손해를 본다. \par
ㄹ. 기회비용의 크기는 비교우위를 결정한다. \par
\end{framedtext}
\begin{choices}
  \item ㄱ, ㄴ
  \item ㄱ, ㄹ
  \item ㄴ, ㄷ
  \item ㄴ, ㄹ
  \item ㄷ, ㄹ
\end{choices}
\end{question}
\answer{2}
